% Ad Atticum 4.9 (ad-atticum-04-09) --- English translation
% Translated by Claude (Anthropic) from the Latin (Perseus).
% Style guide: STYLE.md.

\ciceroLetterOpener{CICERO ATTICO salutem}

\ciceroSection{1} I should very much like to know whether the tribunes are blocking the census by spoiling its days (this is the rumour here), and what is being done and intended about the censorship in general. Here we have been with Pompey. He talked with me at length about politics --- not at all pleased with himself, as he was speaking (for that is how it must be put with him); scornful of Syria, blustering about Spain --- this too, as he was speaking; and indeed I think on every score, when we are talking of him, let it be a kind of \textit{this too} \textit{[Greek: kai tode]}. He thanked you also for taking on the arrangement of the statues; with me he was, by Hercules, most charmingly effusive. He even came over to my villa at Cumae from his own. To me he seemed to want nothing less than for Messalla to stand for the consulship. About this very thing, if you know anything, I should like to know.

\ciceroSection{2} Your writing that you will commend our glory to Lucceius, and our building, which you keep visiting, is welcome. My brother Quintus has written to me that, since you have his sweet little Cicero with you, he will come to you on the Nones of May. I left Cumae on the fifth day before the Kalends of May. That day, at Naples with Paetus. On the fourth day before the Kalends of May, very early in the morning as I was setting out for the Pompeian villa, I wrote this.
